% !TEX root = ../main.tex
\chapter{ケーススタディ：API adapter と自然性}
\label{chap:ch25-api-adapter-naturality}

\begin{chapterabstract}
本章では、APIレスポンス形式の変更を題材に、adapter の正しさを自然性として表す。
旧形式のレスポンスに業務ロジックを適用してから新形式へ変換しても、先に新形式へ変換してから業務ロジックを適用しても、結果が一致する。
この性質を、自然性の四角形と Lean 4 の等式として確認する。
本章の狙いは、adapter を「変換関数があるから安全」と見るのではなく、「業務ロジックと干渉しないことを証明できる層」として扱うことである。
\end{chapterabstract}

\begin{learninggoals}
\item APIレスポンスwrapperの変更を、型に依存しない変換として整理できる。
\item 自然性の四角形を、adapter と \texttt{map} の順序交換として読める。
\item \texttt{OldResponse A} と \texttt{NewResponse A} の小さな Lean モデルを書ける。
\item adapter を通す位置を変えても結果が同じであることを、Lean の定理として表せる。
\item 実務で自然性を使うべき範囲と、使ってはいけない範囲を区別できる。
\end{learninggoals}

\section{APIレスポンス形式の変更}

APIのバージョンアップでは、レスポンスの中身そのものは同じでも、包み方が変わることがある。
たとえば、旧APIでは次のような形で返していたとする。

\begin{itemize}
\item ステータスコードを \texttt{status} に入れる。
\item 実データを \texttt{body} に入れる。
\item トレース用のIDを \texttt{requestId} に入れる。
\end{itemize}

新APIでは、名前や追加メタデータを変えて、次のような形にしたい。

\begin{itemize}
\item ステータスコードを \texttt{code} に入れる。
\item 実データを \texttt{data} に入れる。
\item \texttt{requestId} は引き継ぐ。
\item キャッシュ由来かどうかを \texttt{cached} として追加する。
\end{itemize}

このとき、旧レスポンスを新レスポンスへ変換する adapter を書くこと自体は難しくない。
しかし、本章で確認したいのは「変換関数が存在するか」ではない。
本当に確認したいのは、adapter がレスポンスの中の業務データに余計な影響を与えないかである。

\begin{intuitionbox}
ここでの adapter は、荷物そのものを変える作業ではなく、配送箱を入れ替える作業だと考える。
箱のラベルや形は変わる。
しかし、箱の中身に対して行う業務処理は、箱を入れ替える前に実行しても、後に実行しても、同じ結果になってほしい。
\end{intuitionbox}

\FigurePlaceholder{fig-ch25-response-shape-change.png}{0.92}{旧APIレスポンスから新APIレスポンスへの包み替え}{fig:ch25-response-shape-change}

\begin{softwarebox}
実務上は、次のような場面で同じ問題が現れる。
\begin{itemize}
\item REST API の \texttt{v1} レスポンスを \texttt{v2} レスポンスへ包み替える。
\item SDK の互換層で、旧clientが期待する型を新clientの型へ変換する。
\item GraphQL や gRPC の境界で、内部DTOを公開DTOに包み直す。
\item レスポンスmetadataを追加しつつ、payloadに対する既存ロジックを再利用する。
\end{itemize}
\end{softwarebox}

\section{\texttt{OldResponse A} と \texttt{NewResponse A}}

まず、payload の型を \texttt{A} として抽象化する。
\texttt{A} は、ユーザーでも注文でも商品でもよい。
重要なのは、レスポンスwrapperが payload の型に対して一様に振る舞うことである。

\begin{definitionbox}
\texttt{OldResponse A} は、型 \texttt{A} の値を旧形式で包むレスポンス型である。
\texttt{NewResponse A} は、同じ型 \texttt{A} の値を新形式で包むレスポンス型である。
ここで \texttt{A} を固定しないのは、adapter が特定のpayload型だけでなく、任意のpayload型に対して同じ規則で動くかを見たいからである。
\end{definitionbox}

Leanでは、次のように最小モデルを作る。
ここではHTTPの詳細、JSONエンコード、エラー詳細、認証情報などは扱わない。
自然性を確認するために必要な情報だけを残す。

\begin{listing}[htbp]
\begin{minted}{lean4}
structure OldResponse (A : Type) where
  status : Nat
  body : Option A
  requestId : String

structure NewResponse (A : Type) where
  code : Nat
  data : Option A
  requestId : String
  cached : Bool
\end{minted}
\caption{旧レスポンスと新レスポンスの最小モデル}
\label{lst:ch25-response-types}
\end{listing}

\texttt{OldResponse} と \texttt{NewResponse} は、どちらも \texttt{Type -> Type} という形を持つ。
つまり、payload の型 \texttt{A} を受け取ると、レスポンス型を返す型レベルの関数として読める。
このような形は、圏論では関手を考える入口になる。

\subsection{payloadだけを変換する \texttt{map}}

業務ロジックは、多くの場合、レスポンス全体ではなく payload に対して働く。
たとえば、内部ユーザー型を公開用ユーザー型へ変換する関数は、レスポンスmetadataではなく \texttt{body} や \texttt{data} の中身に作用する。

\begin{listing}[htbp]
\begin{minted}{lean4}
namespace OldResponse

def map {A B : Type} (f : A -> B)
    (r : OldResponse A) : OldResponse B :=
  { status := r.status,
    body := Option.map f r.body,
    requestId := r.requestId }

end OldResponse

namespace NewResponse

def map {A B : Type} (f : A -> B)
    (r : NewResponse A) : NewResponse B :=
  { code := r.code,
    data := Option.map f r.data,
    requestId := r.requestId,
    cached := r.cached }

end NewResponse
\end{minted}
\caption{payloadだけを変換するmap}
\label{lst:ch25-response-map}
\end{listing}

この \texttt{map} は、metadata をそのまま保ち、payload だけに \texttt{f : A -> B} を適用する。
旧レスポンスでも新レスポンスでも、payload が存在しない \texttt{none} の場合は、存在しないままにする。

\begin{leanbox}
ここで \texttt{Option.map f r.body} は、「値があれば \texttt{f} を適用し、なければ \texttt{none} のままにする」という処理である。
レスポンスwrapperの中身にだけ業務ロジックを持ち上げる操作だと読める。
\end{leanbox}

\section{adapter が業務ロジックと可換であること}

次に、旧レスポンスから新レスポンスへの adapter を定義する。
この adapter は、\texttt{status} を \texttt{code} へ、\texttt{body} を \texttt{data} へ写し、\texttt{requestId} を保存する。
新しく追加された \texttt{cached} には、ここでは既定値 \texttt{false} を入れる。

\begin{listing}[htbp]
\begin{minted}{lean4}
def adapt {A : Type} (r : OldResponse A) : NewResponse A :=
  { code := r.status,
    data := r.body,
    requestId := r.requestId,
    cached := false }
\end{minted}
\caption{旧レスポンスから新レスポンスへのadapter}
\label{lst:ch25-adapter}
\end{listing}

ここで重要なのは、\texttt{adapt} が \texttt{A} の中身を見ていないことである。
ユーザーIDが偶数かどうか、注文金額が高いかどうか、文字列が空かどうか、といったpayload固有の判断はしない。
そのため、payload に対する業務ロジックとは独立していると期待できる。

\begin{definitionbox}
本章でいう「可換」とは、二つの処理の順序を入れ替えても最終結果が同じである、という意味である。
ここでは、次の二つの経路が同じであることを示す。
\begin{align*}
&\text{経路1: 旧レスポンス内でpayloadを変換してから、新レスポンスへadaptする。}\\
&\text{経路2: 先に新レスポンスへadaptしてから、新レスポンス内でpayloadを変換する。}
\end{align*}
\end{definitionbox}

payload変換を \(f : A \to B\)、旧レスポンスを \(r : \texttt{OldResponse}\ A\) とすると、証明したい等式は次である。

\[
\mathsf{adapt}\,(\mathsf{OldResponse.map}\ f\ r)
=
\mathsf{NewResponse.map}\ f\,(\mathsf{adapt}\ r).
\]

左辺は「先に旧形式の中で業務ロジックを適用し、その後adapterを通す」経路である。
右辺は「先にadapterを通し、その後新形式の中で業務ロジックを適用する」経路である。

\section{自然性の四角形をコードで読む}

自然変換の自然性は、よく四角形の図で表される。
本章では、この四角形をAPI adapterの検証図として読む。

\FigurePlaceholder{fig-ch25-naturality-square.png}{0.88}{API adapter の自然性の四角形}{fig:ch25-naturality-square}

図の上辺は、旧レスポンス内でpayloadを変換する処理である。
図の下辺は、新レスポンス内でpayloadを変換する処理である。
左辺と右辺は、それぞれpayload型 \texttt{A} と \texttt{B} に対する adapter である。

\begin{intuitionbox}
自然性の四角形は、「中身を変えてから包みを変える」と「包みを変えてから中身を変える」が同じである、という確認図である。
API adapter の文脈では、これは「adapter が業務ロジックの前後どちらに置かれても意味を変えない」という主張になる。
\end{intuitionbox}

\subsection{自然性が言っていること}

自然性は、単に「あるサンプル入力で同じ結果になる」という主張ではない。
任意のpayload型 \texttt{A}、任意の変換先型 \texttt{B}、任意の業務ロジック \(f : A \to B\)、任意の旧レスポンス \(r\) について、同じ等式が成り立つという主張である。

\begin{softwarebox}
テストであれば、\texttt{User} レスポンス、\texttt{Order} レスポンス、\texttt{Product} レスポンスをいくつか用意して確認することになる。
Leanで自然性を証明すると、モデル化した範囲では、特定の型や特定の入力に依存しない一般的な性質として確認できる。
\end{softwarebox}

\subsection{自然性が言っていないこと}

自然性は、adapter が現実のHTTP通信やJSONパーサのすべての振る舞いを保証するわけではない。
本章の証明が扱うのは、型で表したレスポンスwrapperとpayload変換の関係である。

\begin{warningbox}
自然性を証明しても、外部APIのステータスコード仕様、JSONフィールド名の実際の綴り、ネットワーク障害、認証、タイムアウト、ログ出力、副作用の順序までは保証されない。
Leanのモデルが保証するのは、モデルに含めた関数とデータ構造についての性質である。
実装コードとの対応は、別途レビューやテストで確認する必要がある。
\end{warningbox}

\section{Lean で自然性を証明する}

証明方針は単純である。
旧レスポンス \texttt{r} を構造体の各フィールドに分解する。
すると、左辺と右辺はどちらも同じ \texttt{NewResponse} の値になる。
したがって、Leanでは \texttt{rfl} で証明できる。

\begin{listing}[htbp]
\begin{minted}{lean4}
theorem adapt_natural {A B : Type} (f : A -> B)
    (r : OldResponse A) :
    adapt (OldResponse.map f r) =
      NewResponse.map f (adapt r) := by
  cases r
  rfl
\end{minted}
\caption{adapterの自然性}
\label{lst:ch25-adapt-natural}
\end{listing}

この証明は短いが、主張は強い。
\texttt{cases r} によって、\texttt{r} は \texttt{status}、\texttt{body}、\texttt{requestId} というフィールドへ分解される。
分解後、左辺と右辺は同じ構造体リテラルに計算される。
そのため、\texttt{rfl} が使える。

\begin{leanbox}
\texttt{rfl} は「両辺が定義を展開すると同じものになる」場合に使える。
ここでは、adapter と \texttt{map} の定義を展開すると、左右が同じ \texttt{NewResponse} になる。
つまり、自然性は adapter の定義そのものから出てくる。
\end{leanbox}

\subsection{二つの経路に名前を付ける}

実務の説明では、自然性の等式をそのまま見るより、二つの経路に名前を付けたほうが読みやすい。
ここでは、先に業務ロジックを適用する経路を \texttt{routeBefore}、先に adapter を通す経路を \texttt{routeAfter} と呼ぶ。

\begin{listing}[htbp]
\begin{minted}{lean4}
def routeBefore {A B : Type} (f : A -> B)
    (r : OldResponse A) : NewResponse B :=
  adapt (OldResponse.map f r)

def routeAfter {A B : Type} (f : A -> B)
    (r : OldResponse A) : NewResponse B :=
  NewResponse.map f (adapt r)

theorem route_position_independent {A B : Type}
    (f : A -> B) (r : OldResponse A) :
    routeBefore f r = routeAfter f r := by
  exact adapt_natural f r
\end{minted}
\caption{adapterを置く位置が違う二つの経路}
\label{lst:ch25-routes}
\end{listing}

\texttt{route\_position\_independent} は、実務向けには「adapterの位置独立性」と読める。
この定理があると、リファクタリングで adapter の位置を変更しても、payloadだけを見る業務ロジックについては結果が変わらないと説明できる。

\subsection{自然な adapter を構造として持つ}

adapter と自然性証明を一緒に持つ小さな構造を作ることもできる。
この構造は、Mathlib の \texttt{NatTrans} を直接使うものではない。
本章の目的に合わせて、\texttt{OldResponse} と \texttt{NewResponse} の間だけに絞った簡易版である。

\begin{listing}[htbp]
\begin{minted}{lean4}
structure ResponseAdapter where
  app : {A : Type} -> OldResponse A -> NewResponse A
  naturality : {A B : Type} -> (f : A -> B) ->
    (r : OldResponse A) ->
    app (OldResponse.map f r) = NewResponse.map f (app r)

def oldToNewAdapter : ResponseAdapter where
  app := fun {A} r => adapt (A := A) r
  naturality := by
    intro A B f r
    exact adapt_natural f r
\end{minted}
\caption{adapterと自然性証明をまとめる構造}
\label{lst:ch25-response-adapter}
\end{listing}

このようにまとめると、adapter を単なる関数ではなく「可換性の証明を持った互換層」として扱える。
仕様レビューでは、adapter の本体だけでなく、何と可換でなければならないかを明確にできる。

\section{実務での意味：adapter の位置を変えても結果が同じ}

ここまでの定理は、実務では次のように使える。

\begin{itemize}
\item 旧APIレスポンスに対して既存の業務ロジックを適用してから、新APIレスポンスへ変換する。
\item 先に新APIレスポンスへ変換し、新形式の上で同じ業務ロジックを適用する。
\item この二つの経路が一致することを、Leanで確認する。
\end{itemize}

たとえば、内部ユーザー型 \texttt{UserV1} を公開用の \texttt{UserView} へ変換する処理を考える。
この処理はpayloadの中身だけを見て、レスポンスmetadataには触れない。
そのような処理に対して、adapter は前後どちらに置いても同じ結果になる。

\begin{listing}[htbp]
\begin{minted}{lean4}
structure UserV1 where
  id : Nat
  name : String

structure UserView where
  id : Nat

def toUserView (u : UserV1) : UserView :=
  { id := u.id }

theorem user_view_adapter_position_independent
    (r : OldResponse UserV1) :
    routeBefore toUserView r = routeAfter toUserView r := by
  exact route_position_independent toUserView r
\end{minted}
\caption{ユーザーpayload変換に特化した自然性}
\label{lst:ch25-user-view}
\end{listing}

この証明では、\texttt{UserV1} や \texttt{UserView} に固有の難しい議論はしていない。
すでに任意の \texttt{A} と \texttt{B} について証明した自然性を、ユーザー変換に適用しているだけである。
これが抽象化の利点である。

\subsection{設計レビューで確認すべきこと}

自然性を実務に持ち込むときは、次の問いが有効である。

\begin{center}
\begin{tabular}{p{0.31\linewidth}p{0.58\linewidth}}
\toprule
確認する問い & 意味 \\
\midrule
payload型を任意にしてよいか & adapter が特定の型だけを特別扱いしていないかを確認する。 \\
metadataだけを変えているか & 業務データの意味に手を入れていないかを確認する。 \\
\texttt{map} はpayloadだけに作用するか & 業務ロジックの範囲をレスポンスwrapperから分離できているかを確認する。 \\
adapterの追加値は定数か & \texttt{cached := false} のように、payloadに依存しない初期値かを確認する。 \\
外部仕様との差分は明示したか & Leanモデルと実装JSON、HTTP仕様、SDK仕様のズレを文書化する。 \\
\bottomrule
\end{tabular}
\end{center}

\FigurePlaceholder{fig-ch25-adapter-review-boundary.png}{0.90}{Leanで証明する範囲と実務レビューで確認する範囲}{fig:ch25-adapter-review-boundary}

\subsection{自然性が壊れる典型例}

次のような adapter は、自然性を期待しにくい。

\begin{itemize}
\item payload の値に応じて \texttt{code} を変更する。
\item payload の型ごとに別のフィールド変換を行う。
\item \texttt{String} payload だけを特別扱いして整形する。
\item adapter の中でDB参照や外部API呼び出しを行う。
\item 変換前後でpayloadの欠損を補完するため、型 \texttt{A} の既定値を注入する。
\end{itemize}

これらが必ず悪いわけではない。
ただし、その場合は「自然な adapter」としてではなく、payloadに依存する業務変換として扱うべきである。
証明すべき性質も、自然性ではなく、事前条件、事後条件、仕様弱化、あるいは片方向互換性になる。

\begin{warningbox}
実務では、adapter と business logic が混ざりやすい。
自然性を証明したいなら、adapter はできるだけ「包みの変換」に限定する。
payload の意味を変更する処理は、別の関数として切り出し、\texttt{map} で扱うほうが設計が明確になる。
\end{warningbox}

\section{本章のまとめ}

\begin{summarybox}
\begin{itemize}
\item API adapter は、旧レスポンスwrapperから新レスポンスwrapperへの変換としてモデル化できる。
\item \texttt{OldResponse.map} と \texttt{NewResponse.map} は、metadataを保ち、payloadだけに業務ロジックを適用する。
\item 自然性は、\texttt{adapt (OldResponse.map f r) = NewResponse.map f (adapt r)} という等式で表せる。
\item Leanでは、構造体を分解し、定義を展開することで、adapter と \texttt{map} の可換性を証明できる。
\item この証明は、adapter が業務ロジックと干渉しないこと、つまり adapter の位置を変えても結果が同じであることを示す。
\item ただし、外部API仕様、JSONエンコード、ネットワーク、副作用、実装との同期は、Leanモデルの外側で確認する必要がある。
\end{itemize}
\end{summarybox}
